\documentclass[11pt,a4paper]{article}

\usepackage[T1]{fontenc}
\usepackage[utf8]{inputenc}
\usepackage{lmodern}
\usepackage[margin=2.6cm]{geometry}
\usepackage{microtype}

\usepackage{amsmath,amssymb,amsthm}
\usepackage{array}
\usepackage{booktabs}
\usepackage{tabularx}
\usepackage{enumitem}
\usepackage[colorlinks,linkcolor=blue,citecolor=blue,urlcolor=blue]{hyperref}

% Left-aligned expanding column, so tables fill the text block cleanly.
\newcolumntype{L}{>{\raggedright\arraybackslash}X}

\DeclareMathOperator{\Tp}{Tp}
\DeclareMathOperator{\mdeg}{mdeg}
\DeclareMathOperator{\codim}{codim}
\DeclareMathOperator{\Hom}{Hom}
\DeclareMathOperator{\Sym}{Sym}

\newcommand{\Qd}{\mathcal{Q}}
\newcommand{\Nhat}{\widehat{N}}
\newcommand{\eref}{\epsilon_{\mathrm{ref}}}

\theoremstyle{plain}
\newtheorem{proposition}{Proposition}
\newtheorem{lemma}[proposition]{Lemma}
\newtheorem{observation}[proposition]{Observation}
\theoremstyle{definition}
\newtheorem{definition}[proposition]{Definition}
\theoremstyle{remark}
\newtheorem{remark}[proposition]{Remark}

\title{\textbf{Chern++}: computational results on positivity\\[2pt]
       for Morin Thom polynomials}
\author{}
\date{}

\begin{document}
\maketitle
\vspace{-2.4em}

\begin{abstract}
We report on a machine-checked study of the two positivity conjectures for Morin
singularities. The infrastructure computes $\Qd_d$ for $d \le 6$ from the explicit orbit
equations of B\'erczi--Szenes \S7.3, self-checked against every invariant fixed in advance.
On top of it: Rim\'anyi's conjecture is verified for $A_6$ (first time), the two reductions of
the $A_5$ handoff note are shown to survive at $d = 6$ while the proposed prefix-positivity
route is shown to \emph{fail} at $d = 6$, and a notion of \emph{denominator certificate} is
introduced whose successes are proofs and whose failures are theorems. The certificate
machinery reproduces the $d = 4$ strong positivity theorem and Proposition~3 of the handoff
note mechanically. It does not close the $A_6$ tail, and we show this is a limitation of the
technique rather than of the search: after generalising Lemma~1 to an absorption criterion,
$13\,090$ candidate certificates fail. Notation follows \texttt{rimanyi\_positivity.pdf},
\texttt{report.pdf} and \texttt{a5\_weak\_positivity\_handoff.pdf}, cited as [N], [R], [H].
\end{abstract}

\section{Setting and notation}

For the Morin singularity $A_d$, the B\'erczi--Szenes formula expresses $\Tp_{A_d}$ as an
iterated residue whose only $d$-dependent input is $\Qd_d = \mdeg_{T_d}(\mathcal{O}_d)$, the
equivariant multidegree of the Borel orbit closure $\mathcal{O}_d = \overline{B_d \cdot \eref}$
inside $\Nhat_d \subset \Hom(\mathbb{C}^d, \Sym^2 \mathbb{C}^d)$. Normalising $z_d = 1$ and
setting $x_j = z_j/z_{j+1}$, positivity is read off the chamber series
\[
  F_d(x) \;=\;
  \frac{\prod_{m<l}\bigl(1 - z_m/z_l\bigr)\,\Qd_d}
       {\prod_{m+r\le l}\bigl(1 - z_m/z_l - z_r/z_l\bigr)}
  \;=\; \frac{N_d(x)}{\prod_{r}\bigl(1 - f_r(x)\bigr)}
  \;=\; \sum_{\beta \ge 0} A_\beta x^\beta ,
\]
where every $f_r$ has nonnegative coefficients and zero constant term. Write
$\tau(i,j,k,\dots) = (j-i,j,k,\dots)$ for the first adjacent transposition, and $C(M)$ for the
$\ell$-free Chern coefficient of [R, Thm.~5].

\section{Infrastructure}

\subsection{Computing \texorpdfstring{$\Qd_d$}{Q\_d}}

We take the route of [N, \S7.3]: rather than parametrise the orbit and eliminate the group
variables, we write down the quadratic relations of Proposition~7.3 directly, saturate by the
defect-zero coordinates to remove the spurious components that appear from $d = 6$ on
([N, \S7.4]), and read the multidegree off an initial ideal. Runtime for $d \le 6$ is seconds.

The design point worth stating is that every quantity the computation can be checked against is
known in advance, and all of them are checked: multihomogeneity of each relation under the torus
weights; the codimension $\dim \Nhat_d - \binom{d}{2}$ forced by homogeneity of the residue
formula; vanishing on random $B_d$-translates of $\eref$; $\deg \Qd_d$; exactness of the chamber
division; and the normalisation fixing the $c_1^d$ coefficient. The stage refuses to emit an
artifact otherwise. Externally, all eleven classical $A_6$ coefficients are reproduced.

\subsection{Scale of the objects}

Table~\ref{tab:algebra} is the practical reason sampling is no substitute for a proof: roughly
half the coefficients of $N_d$ are negative at every order, and the counts grow fast.

\begin{table}[ht]
\centering
\caption{The mined chamber algebras. $\deg\Qd_d = \dim\Nhat_d - \binom{d}{2}$ is the
codimension of the orbit closure.}
\label{tab:algebra}
\begin{tabularx}{\textwidth}{@{}Lrrrrrrr@{}}
\toprule
 & vars & $\dim\Nhat_d$ & $\deg\Qd_d$ & $\Qd_d$ terms & $N_d$ terms & $\deg N_d$
 & $N_d$ negative \\
\midrule
$A_4$ & 3 & 7  & 1 & 3   & 54    & 13 & 27 \ (50\%) \\
$A_5$ & 4 & 13 & 3 & 19  & 814   & 31 & 414 \ (51\%) \\
$A_6$ & 5 & 22 & 7 & 418 & 26618 & 66 & 13309 \ (50\%) \\
\bottomrule
\end{tabularx}
\end{table}

\subsection{Exactness}

No result below depends on floating point. One point deserves emphasis because it is a trap:
Laurent coefficients are accumulated in \texttt{int64}, and they outgrow it. A wrapped
accumulator reads as a large \emph{negative} number, that is, as a spurious counterexample to
the very conjecture under test. The evaluator therefore replays the fixed point in floating
point as a magnitude probe and raises rather than returning wrapped values; this is what caps
the $A_6$ sweep at $\ell = 5$ below.

\section{Results for \texorpdfstring{$A_6$}{A\_6}}

\begin{proposition}[verified]
Rim\'anyi's conjecture holds for $A_6$ at every relative dimension $\ell \le 5$, and for $A_5$
at $\ell \le 8$.
\end{proposition}

These are finite exact computations, not proofs of the infinite statement. For $A_5$,
[R, Thm.~8] reaches $\ell \le 11$ using unbounded-precision arithmetic; our contribution at
$d = 5$ is speed rather than range, and the $\ell$-free reduction makes the two nested anyway.
The $A_6$ verification is new. Passing $\ell = 5$ requires exact arithmetic beyond
\texttt{int64}; residue arithmetic modulo several primes with CRT reconstruction is the natural
route.

The negative Laurent coefficients of $F_6$ behave much as at $d = 5$; see
Table~\ref{tab:series}. The first counterexample to strong Laurent positivity occurs at $d = 6$
at the same $\beta$ as at $d = 5$, extended by a zero. The last column is the one that matters
in \S\ref{sec:prefix}.

\begin{table}[ht]
\centering
\caption{Negative Laurent coefficients of the chamber series, to total degree $11$.}
\label{tab:series}
\begin{tabularx}{\textwidth}{@{}Lrrlrr@{}}
\toprule
 & terms & negative & first negative & most negative & with $i = 0$ \\
\midrule
$A_4$ & 237 & 0  & ---           & ---  & 0 \\
$A_5$ & 490 & 6  & $(1,1,2,1)$   & $-1$ & 0 \\
$A_6$ & 796 & 10 & $(1,1,2,1,0)$ & $-2$ & 2 \\
\bottomrule
\end{tabularx}
\end{table}

\subsection{The reductions of [H] at \texorpdfstring{$d = 6$}{d = 6}}

\begin{observation}
Both reductions of [H] survive at $d = 6$ throughout the tested range ($\deg \le 11$): the
unpaired tail $i > j \Rightarrow A_{i,j,\dots} \ge 0$, and the paired inequality
$A_{i,j,\dots} + A_{j-i,j,\dots} \ge 0$ for $0 \le i \le j$.
\end{observation}

This is evidence that the $\tau$-pairing frame is structural rather than an $A_5$ accident.

\subsection{Prefix positivity fails at \texorpdfstring{$d = 6$}{d = 6}}
\label{sec:prefix}

[H, \S8] observes that every tested coefficient of $B = F_5/(1-a)$ is nonnegative, and
[H, \S11.4] proposes proving that as a stepping stone, since
$B_{i,j,k,l} = \sum_{r \le i} A_{r,j,k,l}$ turns the paired inequality into a statement about
adjacent differences of a prefix array.

\begin{proposition}
$F_6/(1-a)$ is not coefficientwise nonnegative. Indeed $A_{(0,2,3,2,2)} = -1$.
\end{proposition}

\begin{proof}
The prefix sum runs over $r \le i$, so $B_{0,j,k,l} = A_{0,j,k,l}$ and any negative coefficient
with $i = 0$ survives untouched. $F_5$ has no negative coefficient with $i = 0$ in the tested
range; $F_6$ has two, the first being $A_{(0,2,3,2,2)} = -1$. Confirmed along two independent
code paths, truncated series products in exact integers and the XLA fixed-point grid.
\end{proof}

The route of [H, \S11.4] is therefore specific to $d = 5$. Note that the obstruction is at the
very first level of the prefix sum, so no refinement of the argument recovers it.

\section{Denominator certificates}

\begin{definition}
Let $N / \prod_{r=1}^{R}(1 - f_r)$ with each $f_r$ nonnegative and of zero constant term. An
\emph{order-$k$ certificate} is an identity
\[
  N \;=\; \sum_{|S| \le k} P_S \prod_{r \in S}(1 - f_r),
  \qquad \text{every } P_S \text{ coefficientwise nonnegative,}
\]
over subsets $S \subseteq \{1, \dots, R\}$.
\end{definition}

Dividing through, $F = \sum_S P_S / \prod_{r \notin S}(1 - f_r)$ is a sum of products of
nonnegative series, so $F \ge 0$. Order $0$ is the statement that $N$ itself is nonnegative.
[H, \S10.3] reports a first-order search of this shape coming back infeasible.

\subsection{Why a returned certificate is a proof}

$P_\emptyset$ is not a free unknown: since $\prod_{r \in \emptyset}(1 - f_r) = 1$, it is the
remainder $N - \sum_{S \ne \emptyset} P_S \prod_{r \in S}(1 - f_r)$. We therefore solve a
feasibility LP in the nonempty parts only, snap the solver's floating-point answer to exact
rationals, and recompute $P_\emptyset$ in exact arithmetic. The identity then holds by
construction, and the single remaining question, whether that exact remainder is coefficientwise
nonnegative, is decided exactly. The LP is demoted to a search heuristic; correctness never
depends on it. The output is a finite algebraic identity checkable in any computer algebra
system.

\subsection{Why a failure is a theorem}

\begin{lemma}[projection]
A constraint on a monomial of total degree $T$ involves only the coefficients of the $P_S$ of
degree $\le T$, because every $\prod_{r \in S}(1 - f_r)$ has constant term $1$ and no
negative-degree terms. Truncating both the constraints and the unknowns at degree $T$ is
therefore an exact \emph{projection} of the feasible set, not a further restriction.
\end{lemma}

Consequently infeasibility of the truncated program proves that no certificate of that order
exists at \emph{any} degree. Failed searches become impossibility results.

\subsection{\texorpdfstring{$A_4$}{A\_4}}

\begin{proposition}
$F_4$ admits an order-$4$ certificate with $\deg P_S \le 8$, exhibited explicitly and verified
exactly; and no certificate of order $\le 3$ exists at any degree. The minimal order is $4$.
\end{proposition}

This is an independent computational proof of [N, Thm.~1] by a method using no special feature
of $d = 4$. The lower bound already follows by hand at depth~$1$: writing $m_S$ for the constant
term of $P_S$, matching the constant term of $N = 1 - a - b - c + \cdots$ forces
$\sum_S m_S \le 1$, while matching the coefficients of $a$, $b$, $c$ forces
\[
  \sum_{S \ni f_0} m_S \ge \tfrac12, \qquad
  \sum_{S \ni f_2} m_S \ge 1, \qquad
  \sum_{S \ni f_5} m_S \ge 1,
\]
where $f_0 = 2a$, $f_2 = b + ab$, $f_5 = c + abc$ are the only factors containing the respective
bare monomials. With a budget of $1$ these can only be met if a single $S$ contains all three.
The machine-found certificate saturates all four bounds exactly, its only two constant terms
being $\tfrac12(1-f_0)(1-f_2)(1-f_4)(1-f_5)$ and $\tfrac12(1-f_2)(1-f_4)(1-f_5)(1-f_6)$.

\subsection{\texorpdfstring{$A_5$}{A\_5}: how far the obstruction reaches}

\begin{proposition}
No denominator certificate of order $\le 7$ exists at any degree for $F_5$, nor for
$F_5/(1-a)$.
\end{proposition}

For $F_5$ this is expected, since $F_5 \not\ge 0$, and serves as a control that the engine never
manufactures a certificate for a false statement. For the prefix series, which \emph{is}
coefficientwise nonnegative as far as we compute, it is informative: it extends [H, \S10.3] from
order~$1$ to order~$7$. The bound sharpens with probe depth: depth~$1$ gives $\ge 4$, depth~$2$
gives $\ge 7$, depth~$3$ rules out everything up to $7$.

\section{The unpaired tail}

\subsection{One series}

Set $H(a) = (1-a)/(1-2a)$ and $J_d = F_d / H$. After $H$ is removed, every occurrence of $a$ in
a numerator or denominator monomial is accompanied by $b$: the chamber monomials containing $a$
are $z_1/z_l = ab\cdots$, and only $z_1/z_2 = a$ is bare, which is exactly the factor $1 - 2a$
that $H$ removes. Hence $\deg_a [b^j c^k \cdots] J_d \le j$, and convolving with
$H(a) = 1 + \sum_{n \ge 1} 2^{n-1} a^n$ gives, for $i > j$,
\[
  A_{i,j,k,\dots} \;=\; 2^{\,i-1}\, [b^j c^k \cdots]\, J_d(1/2, b, c, \dots).
\]
The entire unpaired tail, infinitely many coefficients, is therefore one positivity question
about a single series in $d - 2$ variables. This is [H, Lemma~2 and Prop.~3] for $d = 5$, stated
so that it applies verbatim for any $d$.

Our $J_5(1/2,b,c,e)$ agrees with [H, eq.~(9)] coefficient for coefficient to degree $12$, which
is what validates the construction. $J_6(1/2,b,c,d,e)$ is coefficientwise nonnegative as far as
we compute ($\deg \le 10$).

\subsection{Multiplicative certificates}

The two successful proofs in the literature, [N, Thm.~1] and [H, Prop.~3], are
\emph{multiplicative}: products of ratios $(1-u)/(1-v)$, each nonnegative by Lemma~1 because
$v - u \ge 0$. That shape is not a finite sum $\sum_S P_S \prod_{r \in S}(1 - f_r)$ with
$P_S \ge 0$, so the additive certificates of \S4 structurally cannot find them.

The diagnostic is decisive: the additive search is obstructed on $J_5(1/2,\cdot)$, a series
whose nonnegativity \emph{is} Proposition~3. When a method fails on a known theorem, the method
is wrong.

Presenting the numerator in factored form, the Vandermonde part being already a factorisation,
one can search for a matching of numerator factors to denominators satisfying the Lemma~1
condition, cancel whatever divides exactly, and ask whether what remains is nonnegative.

\begin{proposition}
For $d = 5$ this closes: all nine numerator factors pair off, three denominators cancel exactly,
and the leftover numerator is exactly $1$. Proposition~3 is reconstructed mechanically.
\end{proposition}

For $d = 6$ it does not. A maximum matching leaves
\[
  \frac{1 - cde - \tfrac32 bcde}{1 - 2cde},
\]
whose expansion has the negative family $b\,(cde)^{k}$, and no available denominator $v$
dominates $u = cde + \tfrac32 bcde$. Maximum matching is in fact greedy in the wrong way: it
consumes denominators the remainder needs. Returning two ratios to the remainder leaves twelve
Lemma-1 ratios times a remainder ($207$-term numerator over nine denominators) that is
nonnegative as far as we check.

\subsection{Generalising Lemma~1, and a negative result}

Lemma~1 requires the numerator in the shape $1 - u$. The following drops that and lets several
denominators cooperate.

\begin{lemma}[absorption]
Write $N = N_+ - N_-$ with both parts coefficientwise nonnegative, and let $S$ be a set of
denominators with $1 - W = \prod_{v \in S}(1 - v)$. If $N_- \le N_+ \cdot W$ coefficientwise,
then $N / \prod_{\mathrm{all}}(1 - v)$ is coefficientwise nonnegative.
\end{lemma}

\begin{proof}
$N - N_+ \prod_S (1-v) = N_+ W - N_- \ge 0$, so $N = N_+ \prod_S (1-v) + P$ with $P \ge 0$ and
\[
  \frac{N}{\prod_{\mathrm{all}}(1-v)}
  \;=\; \frac{N_+}{\prod_{v \notin S}(1-v)} \;+\; \frac{P}{\prod_{\mathrm{all}}(1-v)},
\]
both nonnegative.
\end{proof}

Lemma~1 is the case $N = 1 - u$, $S = \{v\}$, where the condition reads $u \le v$; the criterion
is strictly stronger, since it applies to numerators not of that shape.

\begin{observation}
Searching jointly over back-off subsets and absorbing subsets, $13\,090$ candidate criteria,
closes nothing for $A_6$.
\end{observation}

Since $J_6(1/2,\cdot)$ is nonnegative as far as we compute, a proof exists. The content of the
observation is that it is not of this form: the technique that proves $d = 5$ does not extend,
and the obstruction is structural rather than a matter of searching harder. Concretely, the
negativity of the residual factor is compensated by the \emph{other} ratios, so no local pairing
argument can succeed.

\section{Assessment}

What the computations suggest about proof strategy:

\begin{enumerate}[itemsep=3pt,topsep=3pt]
\item The $\tau$-pairing reduction is the right frame, both its inequalities surviving at
  $d = 6$, but the prefix-positivity stepping stone is not, and is unavailable from $d = 6$ on.
\item Additive, Positivstellensatz-shaped certificates are obstructed on \emph{every} target
  tried, including provable ones. On present evidence they are not the tool for this problem,
  and we would not recommend further effort extending [H, \S10.3] in that direction.
\item Multiplicative certificates are the tool that works, and they are now mechanised: they
  reproduce both published proofs from scratch. Their reach stops before $A_6$.
\item The natural remaining targets are the paired series $\mathcal{P}$ of [H, eq.~(12)], for
  which the construction is now cheap since $S = J_5(1/2,\cdot)$ is exactly the object above,
  the reflection is an index reflection, and [H, \S7] supplies the denominator; and closed-form
  domination along the negative families of [R, \S7.2], where each family settled is infinitely
  many cases.
\end{enumerate}

\section*{Reproducibility}

Everything above is reproducible from a clean checkout of Chern++. Statements are separated
throughout into finite exact verifications and proofs; the certificate machinery and the
closed-form family arguments are the only components producing the latter. The annotated
walkthrough is \texttt{src/examples.ipynb}, executed with its outputs.

\end{document}
